% ECONOMICS_PROBLEM: {"id": "O31-626", "title": "Cross-disciplinary discovery", "jel_code": "O31", "source_id": "OP-0411"}
% !TeX program = xelatex
\documentclass[11pt,a4paper]{article}
\usepackage[margin=23mm]{geometry}
\usepackage{fontspec}
\setmainfont{Latin Modern Roman}
\usepackage{amsmath,amssymb,mathtools}
\usepackage{xcolor,enumitem,booktabs,longtable,array,xurl,hyperref}
\definecolor{muted}{HTML}{53636C}
\hypersetup{colorlinks=true,urlcolor=blue,linkcolor=blue}
\setlength{\parindent}{0pt}
\setlength{\parskip}{5pt}
\newcommand{\R}{\mathbb R}
\newcommand{\E}{\mathbb E}
\newcommand{\N}{\mathbb N}
\newcommand{\ind}{\mathbf 1}
\newcommand{\Var}{\operatorname{Var}}
\newcommand{\Cov}{\operatorname{Cov}}
\newcommand{\supp}{\operatorname{supp}}
\DeclareMathOperator*{\argmin}{arg\,min}
\DeclareMathOperator*{\argmax}{arg\,max}
\newcommand{\entrymeta}[3]{{\sffamily\small\color{muted}\textbf{\texttt{#1}}\quad|\quad #2\par #3\par}}
\newcommand{\Problemref}[1]{\texttt{#1}}
\newcommand{\JELref}[2]{\texttt{#2} (JEL \texttt{#1})}
\begin{document}
\section*{Cross-disciplinary discovery}
\textbf{Problem ID:} \texttt{O31-626}\quad \textbf{JEL:} \texttt{O31}\par
% BEGIN ECONOMICS STATEMENT
\label{problem:OP-0411}
{\sffamily\small\color{muted}Original subject group: Innovation and AI economics\par}
\label{definitions:problem:30007693}
\textbf{Audit classification:} Explicit published open question. Published cross-disciplinary questions verified. The original monetary-or-scientific endpoint could make its cost comparison dimensionally invalid.
\subsection*{Definitions and precise research target}
Take interdisciplinary research projects funded by a defined grant program, with comparable budgets and a three-year evaluation horizon. Let \(A\in\{0,1\}\) indicate access to a fixed AI research system, and \(M\in\{0,1\}\) indicate assignment to a team combining two prespecified disciplines rather than a team within one discipline. For project \(i\), let \(Y_i(a,m)\) be the monetary or preregistered scientific value of independently validated discoveries under assignment \((a,m)\), including zero for projects producing none. Define the cross-disciplinary complementarity
\[
\Delta=E[Y_i(1,1)-Y_i(0,1)-Y_i(1,0)+Y_i(0,0)].
\]
The expectation concerns projects eligible for the program; randomized team composition and AI access would identify this finite-population interaction under consistency and negligible spillovers between assigned teams. Where ideas diffuse, define connected team clusters and randomize at that level instead. Measure disciplinary distance using a fixed pre-intervention taxonomy, and evaluate novelty using blinded reviewers rather than citations generated by the intervention itself. Let \(K(m)\) denote the additional coordination cost of mixed teams. Determine whether the distribution of validated gains and \(\Delta\) justify \(K(1)-K(0)\), and which institutional allocation of expertise supports useful discoveries. Positive interaction in this setting would not establish that AI improves every interdisciplinary project.

\textbf{Definitions, assumptions and mathematical qualifications.} Choose either monetary discovery value or one specified scientific index with fixed units; do not interchange them within \(\Delta\). Team assignment is a reproducible policy drawing members from prespecified disciplines and matching project budget, size and eligibility. \(Y_i(a,m)\) includes failed projects and has finite mean. For a fixed cohort of \(n\) projects, \(E[Y_i(a,m)]\) denotes \(n^{-1}\sum_iY_i(a,m)\); superpopulation claims require a separate sampling law. To compare coordination costs measured in dollars with discovery outcomes, value outcomes in dollars or specify the willingness-to-pay conversion. Define net value \(Y_i(a,m)-C_i(a,m)\) separately from the gross interaction.
\subsection*{Variants and assumption changes}
\begin{enumerate}
\item \textbf{Editorial, scientific endpoint.} Use blinded reproducibility and a frozen disciplinary taxonomy to identify the scientific-index interaction; report coordination resources separately, without a money-benefit comparison.
\item \textbf{Editorial, net value and spillovers.} Use monetized validated outcomes and net costs, randomize connected research clusters, and compare the AI-by-mixed-team interaction in net value rather than citations.
\end{enumerate}
\subsection*{References and source audit}
\textbf{Checked source.} 2024 Stanford draft, Section 3.2, PDF pp. 10--11. Full primary draft and cover retrieved. \url{https://digitaleconomy.stanford.edu/app/uploads/2024/11/ETAI-White-Paper.pdf}.
\begin{enumerate}\raggedright
\item\hypertarget{definitions:source:30007693:1}{} Erik Brynjolfsson; Anton Korinek; Ajay Agrawal. 2024. \emph{The Economics of Transformative AI: A Research Agenda}. 2024 Stanford draft, Section 3.2, PDF pp. 10--11.
\url{https://digitaleconomy.stanford.edu/app/uploads/2024/11/ETAI-White-Paper.pdf}.
\end{enumerate}

\subsection*{Common conventions: Source mathematical conventions}

These conventions apply to each entry unless explicitly replaced.

\textbf{Models and identification.} A primitive model \(m\) specifies all probability laws, feasible choices, information, equilibrium and measurement rules used by its target. The admissible class \(\mathcal M\) is fixed independently of outcomes. If \(P_O(m)\) is its observable probability law and \(\tau(m)\) the target, its identified set at observable law \(P\) is
\[
 \mathcal I_\tau(P)=\{\tau(m):m\in\mathcal M,\ P_O(m)=P\}.
\]
Point identification means that this set is a singleton. An empty set signals incompatibility of data and assumptions. Coordinate bounds need not describe the joint set. Bounds are sharp only if every retained target is attainable or arbitrarily approachable by compatible models. Finite bounds require explicit restrictions on unobserved quantities; unrestricted confounding, hidden wealth, extrapolation or arbitrary equilibrium selection can yield uninformative sets.

\textbf{Causal interventions.} A potential outcome \(Y_i(p)\) is the outcome under a complete policy regime \(p\), including timing, financing, information and market spillovers. The target population and units are fixed before treatment. With interference, use the complete assignment vector or a justified exposure mapping; individual no-interference notation alone cannot identify an economy-wide intervention. A difference of mean potential outcomes does not require the cross-world joint distribution. Conditioning on post-treatment migration, survival, enrollment or employment changes the estimand and needs separate justification.

\textbf{Statistical statements.} An estimator, confidence set, forecast or test specifies a sampling law, model class and loss/coverage criterion. Uniform \((1-\alpha)\)-coverage means \(\inf_{m\in\mathcal M}\Pr_m\{\tau(m)\in C_n\}\geq1-\alpha\), or an explicitly asymptotic version. The whole parameter space trivially has coverage, so informativeness requires a stated width, risk or power objective. A forecasting improvement is relative to a named benchmark, information set and evaluation population.

\textbf{Equilibrium and optimization.} An equilibrium comprises feasible plans, optimality, beliefs where needed, budget constraints and all market/resource clearing. The primitives or a stated benchmark define each correspondence \(\mathcal E(p;m)\). Nonemptiness is assumed or proved, never inferred just from compact policy sets. A selected-equilibrium target uses a declared selection rule; a robust target quantifies over all permitted equilibria. Use suprema or infima unless attainment is established. An \(\varepsilon\)-optimal policy is feasible and within \(\varepsilon\) of the finite optimum value. Report an empty feasible set separately.

\textbf{Welfare and accounting.} Normative utility, interpersonal normalization, social weights, numeraire, discounting and the use of public receipts are primitives. Behavioral utility need not coincide with the welfare criterion. Household welfare includes ownership income and rents once; adding profits again to compensating variation generally double-counts them. A monetary equivalent is defined by an explicit transfer and price convention, with monotonicity and range conditions. Average effects, mechanism contributions, transfers and welfare losses are different targets. Infinite sums require convergence and dynamic feasible plans require terminal or no-Ponzi conditions.

\textbf{Source versions.} A working-paper date, revision date and journal publication year are distinguished. A DOI for a journal version is not automatically the DOI of an earlier manuscript. Locators refer to the stated version and printed pages where specified. A metadata-only check does not verify a claim about a full-text conclusion. Every entry's source subsection narrows the provenance claim accordingly.
% END ECONOMICS STATEMENT
\end{document}
